\section*{Sets and Indices}

Let
\[
M = \{A,B,C\}
\]
be the set of manufacturers.

No additional indexing sets are used in the implementation.

\section*{Parameters}

For each manufacturer \(m \in M\):
\begin{itemize}
    \item \(c_m\): cost per chair.
    \item \(q_m\): chairs per order.
    \item \(U_m\): upper bound on the number of orders, computed in code as
    \[
    U_m = \left\lfloor \frac{\texttt{max\_total}}{q_m} \right\rfloor + 1.
    \]
\end{itemize}

Using the provided data:
\[
c_A = 50,\quad c_B = 45,\quad c_C = 40,
\]
\[
q_A = 15,\quad q_B = 10,\quad q_C = 10.
\]

Global parameters:
\begin{itemize}
    \item \(L\): minimum total chairs required (\texttt{min\_total\_chairs}), with \(L=100\).
    \item \(U\): maximum total chairs allowed (\texttt{max\_total\_chairs}), with \(U=500\).
    \item \(r_{BA}\): minimum number of chairs from manufacturer \(B\) if manufacturer \(A\) is used (\texttt{min\_chairs\_B\_if\_A}), with \(r_{BA}=10\).
    \item \(M^{\text{big}}\): big-\(M\) constant used for linking, defined in code as
    \[
    M^{\text{big}} = U + 1 = 501.
    \]
\end{itemize}

The order upper bounds used by the code are therefore
\[
U_A = \left\lfloor \frac{500}{15} \right\rfloor + 1 = 34,\qquad
U_B = \left\lfloor \frac{500}{10} \right\rfloor + 1 = 51,\qquad
U_C = \left\lfloor \frac{500}{10} \right\rfloor + 1 = 51.
\]

\section*{Decision Variables}

\begin{itemize}
    \item \(x_A\) (code: \texttt{x\_A}), \(x_B\) (code: \texttt{x\_B}), \(x_C\) (code: \texttt{x\_C}): integer number of orders placed with manufacturers \(A,B,C\), respectively.
    \[
    x_m \in \mathbb{Z}_{\ge 0}, \qquad x_m \le U_m \quad \forall m \in M.
    \]
    \item \(y_A\) (code: \texttt{y\_A}), \(y_B\) (code: \texttt{y\_B}), \(y_C\) (code: \texttt{y\_C}): binary indicators equal to 1 if the corresponding manufacturer is used.
    \[
    y_m \in \{0,1\} \quad \forall m \in M.
    \]
\end{itemize}

The code also defines implied chair quantities:
\[
\text{chairs}_A = q_A x_A,\qquad \text{chairs}_B = q_B x_B,\qquad \text{chairs}_C = q_C x_C.
\]

\section*{Objective}

Minimize total purchasing cost:
\[
\min \; c_A q_A x_A + c_B q_B x_B + c_C q_C x_C.
\]

With the given data, this is
\[
\min \; 50\cdot 15\,x_A + 45\cdot 10\,x_B + 40\cdot 10\,x_C.
\]

\section*{Constraints}

Total chairs must lie within the required interval:
\[
q_A x_A + q_B x_B + q_C x_C \ge L,
\]
\[
q_A x_A + q_B x_B + q_C x_C \le U.
\]

Binary--integer linking constraints used in the code:
\[
x_A \le M^{\text{big}} y_A,\qquad x_A \ge y_A,
\]
\[
x_B \le M^{\text{big}} y_B,\qquad x_B \ge y_B,
\]
\[
x_C \le M^{\text{big}} y_C,\qquad x_C \ge y_C.
\]

Dependency from \(A\) to \(B\): if manufacturer \(A\) is used, then at least \(r_{BA}\) chairs must be ordered from \(B\):
\[
q_B x_B \ge r_{BA}\, y_A.
\]

Dependency from \(B\) to \(C\): if manufacturer \(B\) is used, then manufacturer \(C\) must also be used:
\[
y_B \le y_C.
\]

Variable domains:
\[
x_A \in \mathbb{Z}_{\ge 0},\; x_B \in \mathbb{Z}_{\ge 0},\; x_C \in \mathbb{Z}_{\ge 0},
\]
\[
x_A \le U_A,\; x_B \le U_B,\; x_C \le U_C,
\]
\[
y_A,y_B,y_C \in \{0,1\}.
\]

\section*{Notes on Implementation}

The implementation builds this model as a mixed-integer linear program in PuLP and solves it with Gurobi via \texttt{GUROBI\_CMD}. The linking formulation is exactly as coded: both
\[
x_m \le M^{\text{big}} y_m
\quad\text{and}\quad
x_m \ge y_m
\]
are imposed for each manufacturer \(m\), so \(y_m=1\) forces at least one order, and \(x_m=0\) forces \(y_m=0\).

The chair quantities \(\text{chairs}_A,\text{chairs}_B,\text{chairs}_C\) are not separate optimization variables in the code; they are linear expressions \(q_m x_m\). The objective uses cost per chair multiplied by these expressions, so the effective cost per order is \(c_m q_m\).

The upper bounds \(U_A,U_B,U_C\) are taken directly from the code logic
\[
\left\lfloor \frac{U}{q_m} \right\rfloor + 1,
\]
even though the total-chair bound already limits feasible orders.
